\subsection{正则空间上的等价关系}

命题 14. 设 X 是正则空间，R 是 X 上的一个闭等价关系。则 R 在 X × X 中的图象 C 是闭的。

设 \((a, b)\) 是 \(X \times X\) 中属于 C 的闭包的一点，并设 V（相应地，W）是 \(a\) 的一个闭邻域（相应地，\(b\) 的一个邻域）；则存在一点 \((x, y) \in C \cap (V \times W)\)。既然 \(x \in V\)，\(y\) 就属于 V 关于 R 的饱和集 S；因此 \(b\) 的每个邻域 W 都与 S 相交。由假设 S 是闭的，从而 \(b \in S\)。现设 B 是 \(\{b\}\) 关于 R 的饱和集，则 \(a\) 的每个闭邻域 V 都与 B 相交；由假设 B 是闭的且 X 是正则的，由此得出 \(a \in B\)，从而 \((a, b) \in C\)。证明完毕。

推论. 在正则空间上，每个既开又闭的等价关系都是豪斯多夫的。

这由命题 14 及第 3 目命题 8 得出。

命题 15. 设 X 是正则空间，F 是 X 的一个闭子集，R 是把 F 的所有点等同起来而得到的 X 上的等价关系{[}换句话说，这个等价关系的等价类是 F（若 \(\mathbf{F} \neq \emptyset\)）以及诸集合 \(\{x\}\)，其中 \(x \in [\mathbb{C}F]\) 。则商空间 X/R 是豪斯多夫的。

设 M 与 N 是 X 中两个不同的等价类。若它们各由 F 的余集中的单点组成，则在豪斯多夫子空间 \(\complement F\) 中存在 M 与 N 的两个不相交的开邻域；它们是 M 与 N 在 X 中的邻域，且关于 R 饱和。若 M = F（于是 \(F \neq \varnothing\)）而 \(N = \{b\}\)，其中 \(b \notin F\)，则既然 X 是正则的，就存在 b 的一个开邻域与 F 的一个开邻域，二者不相交；这些邻域关于 R 饱和，命题得证。

注意，商空间 X/R 未必是正则的（第九章，§ 4，习题 14）。

